\documentclass[11pt, a4paper]{article}
\usepackage[utf8]{inputenc}
\usepackage[spanish]{babel}
% Ajustamos bottom a 3.5cm para darle espacio al pie de página con la imagen
\usepackage[left=2.5cm, right=2.5cm, top=2.5cm, bottom=3.5cm]{geometry} 
\usepackage{graphicx} % Necesario para insertar la imagen de la firma y del pie
\usepackage{tabularx}
\usepackage{array}
\usepackage{parskip}
\usepackage{fancyhdr} % Paquete para encabezados y pies de página

% ==========================================
% CONFIGURACIÓN DEL PIE DE PÁGINA
% ==========================================
\pagestyle{fancy}
\fancyhf{} % Limpia encabezados y pies por defecto

% Elimina las líneas horizontales de encabezado y pie
\renewcommand{\headrulewidth}{0pt}
\renewcommand{\footrulewidth}{0pt}

% Firma / Imagen al lado derecho del pie de página
\fancyfoot[R]{\FirmaPie}
% Mantiene el número de página al centro (opcional)
\fancyfoot[C]{\thepage}

% ==========================================
% DEFINICIÓN DE VARIABLES
% ==========================================
%%INICIO-VARIABLES%%
%%FIN-VARIABLES%%
% ==========================================

\begin{document}

% --- TÍTULO ---
\begin{center}
    \textbf{\Large INFORME \NumeroInforme}
\end{center}

\vspace{2em}

% --- ENCABEZADO FORMATO MEMORÁNDUM / INFORME ---
\renewcommand{\arraystretch}{1.6} % Controla el espaciado vertical entre filas
\begin{tabularx}{\textwidth}{@{} p{3.5cm} c X @{}}
    \textbf{Para}       & : & \ParaNombre \\
                        &   & \ParaCargo \\[0.2em]
    \textbf{De}         & : & \DeNombre \\[0.2em]
    \textbf{Asunto}     & : & \AsuntoTexto \\[0.2em]
    \textbf{Referencia} & : & \ReferenciaTexto \\[0.2em]
    \textbf{Fecha}      & : & \FechaTexto \\
\end{tabularx}

\vspace{1.2em}
% Línea divisoria horizontal
\hrule height 0.6pt
\vspace{1.5em}

% --- CUERPO DEL INFORME ---

El presente informe tiene como objetivo describir las actividades desarrolladas en concordancia con lo establecido en los Términos de Referencia, detallando las acciones ejecutadas durante el período correspondiente.


\section{Antecedentes}

\begin{itemize}
    \item El Instituto Tecnológico de la Producción – ITP viene ejecutando el proyecto denominado “Optimización de la gestión productiva y comercial mediante la adopción de tecnologías digitales y sostenibles para incrementar la competitividad de las unidades productivas de la cadena de valor de calzado y marroquinería en el ámbito del CITECCAL Arequipa”, en el marco del Programa de Fortalecimiento Estratégico de Servicios Tecnológicos – FEST.
    \item El referido proyecto contempla la ejecución de servicios tecnológicos orientados al fortalecimiento de la gestión productiva, diseño y modelaje, soporte productivo, digitalización comercial y sostenibilidad para unidades productivas de la cadena de cuero, calzado y marroquinería.
    \item Asimismo, dentro de las acciones del proyecto se considera la ejecución de servicios de diseño y modelaje de productos innovadores de marroquinería incorporando cuero recuperado, así como servicios de soporte en producción para nuevos e innovadores productos básicos de marroquinería, en el marco de estrategias de sostenibilidad y economía circular.
    \item En ese contexto, se desarrollaron actividades técnicas vinculadas al diseño, modelaje, patronaje, seriado y prototipado de productos de marroquinería y calzado, incorporando tecnologías digitales, optimización de procesos y aprovechamiento de materiales recuperados, conforme a los objetivos y actividades establecidos en el proyecto FEST.
\end{itemize}


\subsection{Actividades a realizar:}
El proveedor del servicio elaborará las siguientes actividades:
\begin{enumerate}
%%INICIO-LISTA_ACTIVIDADES%%
%%FIN-LISTA_ACTIVIDADES%%
\end{enumerate}

\subsection{Plazo de Ejecución:}

El plazo máximo de ejecución del servicio será hasta los \PlazoEjecucion\ días calendario, contabilizados a partir del día siguiente hábil de suscrito el contrato.


\subsection{\EntregableTexto\ Entregable:}

Hasta los \PlazoEntregable\ días calendario, contabilizados a partir del día siguiente hábil de suscrito el contrato.

\section{Análisis}

De acuerdo con lo señalado en los antecedentes, corresponde precisar que la suscripción del contrato entre la Entidad y el contratista origina una relación jurídica patrimonial, mediante la cual el contratista asume la obligación de ejecutar la prestación contratada a favor de la Entidad; mientras que esta última asume la obligación de cumplir con la contraprestación correspondiente, consistente principalmente en la retribución económica por los servicios efectivamente prestados conforme a los términos establecidos en el contrato y los TDR.

En ese sentido, en cumplimiento del contrato se desarrollaron las siguientes actividades:

%%INICIO-ACTIVIDADES%%
%%FIN-ACTIVIDADES%%

\section{CONCLUSIONES:}

Que, habiendo desarrollado las actividades antes mencionadas, el suscrito ha cumplido con la presentación en estricto cumplimiento de lo establecido en los términos de referencia del presente servicio, que forma parte de la Orden de Servicio N° \os\ SIAF: \siaf.\\

Sin otro particular, quedo de usted a la espera de su consideración y posterior conformidad del servicio, para que se prosiga con el trámite que corresponde.\\


% --- BLOQUE DE FIRMA ---
\begin{flushleft}
    \FirmaBloque

    Atentamente, \\[0.3em]
    \DeNombre \\
    \textbf{DNI:} \MiDNI \\
    \textbf{RUC:} \MiRUC
\end{flushleft}

\vspace{1.5em}

Se adjunta:

\begin{enumerate}
    \item Orden de servicio
    \item Recibo por honorarios
    \item RNP
    \item Autorización de CCI
    \item Ficha RUC
    \item Suspensión de 4ta
    \item TDR
\end{enumerate}

\end{document}